% docs/.paper/abstract.tex
\begin{abstract}
Repository-scale software translation across programming languages is a necessary task in software engineering. While Large Language Models (LLMs) translate isolated functions accurately, migrating multi-file repositories introduces cross-file type dependencies, cyclic imports, and context saturation that defeat single-shot prompting. Existing multi-agent frameworks close this gap with closed-source cloud models---incurring monetary cost, execution latency, and data exposure---or they restrict themselves to a single source--target language pair because adapting each new pair requires a large bespoke engineering investment.

We present \textbf{MAgHARCM} (\textbf{M}ulti-\textbf{Ag}ent \textbf{H}arness for \textbf{A}rchaeology-guided \textbf{R}epository-scale \textbf{C}ode \textbf{M}igration). MAgHARCM is an open-source translation harness written in Go that executes entirely on local Small Language Models (SLMs) served by Ollama. It coordinates four specialised agents---Analyzer, Planning, Translator, and Validator---through an explicit \textsc{Analyze} $\to$ \textsc{Plan} $\to$ \textsc{Translate} $\to$ \textsc{Validate} $\to$ \textsc{Repair} loop. A reverse-topological dependency DAG with back-edge cycle removal enforces callee-before-caller ordering, a $\le 4$\,KB Symbol Navigator and a $\le 4$\,KB Prior-Modules Memory cap context per chunk to prevent saturation on small coding models, and a Validator-driven repair cascade combines compiler diagnostics, an adversarial test-weakening guard, and a CodaMOSA-style coverage plateau detector.

We evaluate MAgHARCM on four open-source benchmark repositories spanning C, Go, and Java targeting Rust across $K{=}3$ stochastic runs. Local pairings of a 30B reasoning model (\texttt{qwen3:30b-thinking}) with an instruction-tuned 4B coding model (\texttt{Qwen3-4B-Instruct-UD-Q4\_K\_XL}) achieve clean compilation on every project, a 100\% test pass rate on algorithmic code (GildedRose), 62.5\% on streaming libraries (Gohistogram), 57.1\% on statistical suites (Stats), and 26.5\% on enterprise validation frameworks (Apache Commons Validator) before the plateau detector terminates the repair loop. We formalise a failure taxonomy identifying borrow-checker lifetime errors (38\%) and cyclic package references (26\%) as the dominant bottlenecks for local-SLM repository migration.
\end{abstract}
